\documentclass[11pt,letterpaper]{article}
% HSE 711 Week 4 — student-editable LaTeX learning and lab template
% Compile: pdflatex -interaction=nonstopmode Week_4_Bash_LaTeX_Template.tex
% Compile twice for stable references. Requires standard TeX packages below.
% Never place participant-level sensitive data in a public report.
\usepackage[margin=1in]{geometry}
\usepackage[T1]{fontenc}
\usepackage[utf8]{inputenc}
\usepackage{lmodern}
\usepackage{microtype}
\usepackage{amsmath,amssymb}
\usepackage{booktabs,longtable,array}
\usepackage{xcolor}
\usepackage{hyperref}
\usepackage{fancyvrb}
\usepackage{listings}
\usepackage{graphicx}
\usepackage{enumitem}
\hypersetup{colorlinks=true,urlcolor=blue,linkcolor=blue}
\definecolor{codebg}{RGB}{247,248,250}
\lstdefinestyle{hsebash}{
  language=bash,basicstyle=\ttfamily\small,
  backgroundcolor=\color{codebg},frame=single,
  breaklines=true,showstringspaces=false,
  columns=fullflexible,keepspaces=true,
  commentstyle=\color{gray}
}
\lstset{style=hsebash}
\newcommand{\student}{[Student name]}
\newcommand{\course}{HSE 711: Foundations in Data Science}
\newcommand{\reportdate}{[Submission date]}
\newcommand{\fillin}[1]{\textbf{[Insert #1]}}
\title{Week 4: Introduction to Bash\\\large Guided Learning Notes and Reproducible Metadata Lab}
\author{\student\\\course}
\date{\reportdate}
\begin{document}
\maketitle
\begin{abstract}
This student-editable report follows the Week 4 lecture on the Unix shell and
Bash, then walks through five exercises using the synthetic
\texttt{pseudo\_metadata.csv} file. It separates commands, expected behavior,
observed results, and methodological interpretation. Replace all bracketed
placeholders only after running commands on your own copy of the data.
\end{abstract}
\tableofcontents
\newpage

\section{Learning objectives}
By the end of this exercise, the student should be able to:
\begin{enumerate}
\item Explain the purpose of a shell and the role of Bash in reproducible data processing.
\item Navigate the filesystem and inspect text data without opening an entire large file.
\item Use quoting, variables, conditionals, and pipelines safely.
\item Subset synthetic metadata and document the transformation.
\item Explain why a file line count is not necessarily a participant count.
\item Distinguish simple delimiter processing from full CSV parsing.
\item Preserve a record of commands, assumptions, outputs, and cleanup operations.
\end{enumerate}

\section{Lecture context and conceptual foundation}
A shell interprets commands and interacts with the operating system. The Week 4
lecture introduces Bash as an interactive and scriptable shell and situates it
historically in the lineage of the Bourne, C, and Korn shells. Its scientific
purpose is to orchestrate file operations and reproducible workflows, including
launching R analyses and computational jobs. Bash does not replace a validated
statistical model in R.

\subsection{Command map}
\begin{longtable}{p{.18\linewidth}p{.35\linewidth}p{.35\linewidth}}
\toprule
Command & Function & Scientific use\\ \midrule
\endhead
\texttt{pwd}, \texttt{cd} & Identify or change working directory & Prevent path confusion\\
\texttt{ls -lah} & List files and metadata & Inspect source/output directories\\
\texttt{head}, \texttt{tail}, \texttt{less} & Preview file contents & Assess structure without full import\\
\texttt{wc -l} & Count newline characters & Initial record-count check\\
\texttt{mkdir}, \texttt{cp}, \texttt{mv} & Organize files & Maintain clear provenance\\
\texttt{awk}, \texttt{cut} & Select rows or fields & Inspect simple tabular data\\
\texttt{grep} & Match text patterns & Search sample or gene labels\\
\texttt{Rscript} & Launch R code & Separate orchestration from statistics\\
\bottomrule
\end{longtable}

\section{Environment and source-data assumptions}
\textbf{Required input:} the course-provided \texttt{pseudo\_metadata.csv},
described as synthetic metadata for ten participants. Its exact column headers,
values, and category coding must be verified before filtering. The original
exercise refers to both \texttt{only\_female.txt} and
\texttt{females\_metadata.csv}; this template explicitly creates the latter
from the former.

\noindent\textbf{Important limitation:} the AWK and CUT demonstrations below
assume simple comma-separated rows without quoted commas or embedded newlines.
A CSV-aware parser should be used if the actual file does not meet that
assumption.

\begin{lstlisting}
bash --version
pwd
ls -lah
head -n 3 pseudo_metadata.csv
\end{lstlisting}
\noindent\fillin{operating system, Bash version, and file header}

\section{Exercise 1: Count lines}
The operation \texttt{wc -l} counts line endings. For a standard CSV with one
header and ten one-line participant records, eleven lines would be expected,
but that is not a measured result until the command is run.
\begin{lstlisting}
wc -l pseudo_metadata.csv
\end{lstlisting}
\noindent\textbf{Observed command output:} \fillin{line count}\\
\textbf{Interpretation:} \fillin{number of data rows and header treatment}

\section{Exercise 2: Subset female participants}
This example locates a column titled \texttt{Sex}, compares labels
case-insensitively, and retains the header. Confirm the actual header first.
\begin{lstlisting}
awk -F ',' '
 NR == 1 {
   for (i=1; i<=NF; i++) {
     h=$i; gsub(/\r/, "", h)
     if (tolower(h)=="sex") sex_col=i
   }
   if (!sex_col) { print "Sex column missing" > "/dev/stderr"; exit 2 }
   print; next
 }
 {
   value=$sex_col; gsub(/\r/, "", value)
   if (tolower(value)=="female") print
 }
' pseudo_metadata.csv > only_female.txt
head only_female.txt
wc -l only_female.txt
\end{lstlisting}
\noindent\textbf{Observed selected rows:} \fillin{female participants}\\
\textbf{Validation:} \fillin{header, category labels, missingness}\\
\textbf{Interpretation:} A filtered file is a selection of records, not a
change in the original participants.

\section{Exercise 3: Create a documented working directory}
\begin{lstlisting}
cp only_female.txt females_metadata.csv
mkdir -p new_dir
cp females_metadata.csv new_dir/
cat > new_dir/notes.txt <<'NOTES'
Source: pseudo_metadata.csv, Week 4 synthetic exercise
Filter: Sex == Female (case-insensitive)
Output: females_metadata.csv
Review: verify header, total rows, and missing categories
NOTES
\end{lstlisting}
\noindent\textbf{Notes review:} \fillin{provenance and filtering assumptions}

\section{Exercise 4: Extract the third and fourth columns}
Column positions have no inherent scientific meaning. Verify their names.
\begin{lstlisting}
cd new_dir
head -n 1 females_metadata.csv
cut -d ',' -f 3,4 females_metadata.csv > females_metadata_sub.csv
head -n 5 females_metadata_sub.csv
\end{lstlisting}
\noindent\textbf{Column names:} \fillin{third and fourth header fields}\\
\textbf{Interpretation:} \fillin{meaning of extracted columns}

\section{Exercise 5: Inspect and safely clean up}
\begin{lstlisting}
ls -lah
cd ..
pwd
ls -lah new_dir
# Deliberately interactive; review each deletion:
rm -ri new_dir
\end{lstlisting}
\noindent\textbf{Cleanup record:} \fillin{which files were removed and retained}\\
\textbf{Safety discussion:} Relative paths and recursive deletions can
destroy unrelated files when the working directory is misunderstood.

\section{Bash as an R workflow launcher}
\begin{lstlisting}
Rscript --vanilla analysis.R input.csv output.csv 10
\end{lstlisting}
A downstream R script can read command-line arguments using
\texttt{commandArgs(trailingOnly = TRUE)}. In this learning exercise, that
command is illustrative: do not claim that an analysis ran unless the R file
and its input are actually available.

\section{Results and evidence table}
\begin{tabular}{p{.23\linewidth}p{.37\linewidth}p{.30\linewidth}}
\toprule
Task & Recorded result & Interpretation / concern\\ \midrule
1: Count & \fillin{lines} & \fillin{header and records}\\
2: Subset & \fillin{selected rows} & \fillin{selection validation}\\
3: Copy & \fillin{files created} & \fillin{provenance}\\
4: Columns & \fillin{field names} & \fillin{meaning and limits}\\
5: Cleanup & \fillin{remaining files} & \fillin{safety}\\
\bottomrule
\end{tabular}

\section{Short scholarly narrative}
\noindent\textit{Ready-to-complete paragraph:}

\begin{quote}
The Week 4 exercise used a synthetic participant metadata table to demonstrate
a reproducible command-line workflow. Initial inspection identified
\fillin{verified file structure}. Filtering selected \fillin{count} records
matching the documented female category, with \fillin{missingness findings}.
The resulting subset was copied to a working directory, documented, and reduced
to the third and fourth source columns, identified as \fillin{field names}.
The exercise illustrates why record selection, field provenance, and deliberate
cleanup are necessary even in a small study. Because these data were synthetic,
no empirical clinical or population-level inference is warranted.
\end{quote}

\section{Reproducibility and submission checklist}
\begin{itemize}[label=$\square$]
\item Input file and true headers verified; no sensitive records published.
\item Bash version, working directory, and commands recorded.
\item Header retained; row counts and selection values inspected.
\item Output names reconcile the two filenames in the exercise prompt.
\item Column positions validated against actual field names.
\item Cleanup performed only after verifying the disposable target directory.
\item All numeric results above reflect actual execution, not expectations.
\item PDF inspected for code clipping, page overflow, and legible tables.
\end{itemize}

\section{References and further study}
\begin{enumerate}
\item HSE 711, Week 4 lecture: \emph{Introduction to Bash Scripting}
(supplied course material).
\item HSE 711, Week 4: \emph{Bash Programming In Class Exercises}
(supplied course material).
\item GNU Project. \emph{Bash Reference Manual}.
\url{https://www.gnu.org/software/bash/manual/bash.html}
\item Software Carpentry. \emph{The Unix Shell}.
\url{https://swcarpentry.github.io/shell-novice/}
\end{enumerate}

\end{document}
